\cleardoublepage{}
\sectionnonum{本科生毕业论文（设计）任务书}

{
    \bfseries
    \noindent 一、题目：面向科学场景的智能体工具调用优化研究\\[1ex]
    \noindent 二、指导教师对毕业论文（设计）的进度安排及任务要求：

    \mdseries\normalsize

    \textbf{进度安排：}

    \vspace{0.5ex}
    \small
    \begin{tabular}{p{0.18\linewidth}p{0.70\linewidth}}
        \toprule
        时间 & 工作任务 \\
        \midrule
        第 1--2 周  & 文献调研与技术选型，明确研究边界，完成开题报告 \\[0.3ex]
        第 3--7 周  & 设计并实现标准化科学工具库，封装核心计算工具，设计任务拆解 Prompt 模板 \\[0.3ex]
        第 8--11 周 & 实现检索增强工具选择算法（RATS）与结构化 DAG 任务规划器，完成系统集成联调与对比实验 \\[0.3ex]
        第 12--14 周 & 构建测试数据集，开展真实场景验证，整理实验数据与分析结果 \\[0.3ex]
        第 15--16 周 & 撰写并修改毕业论文，整理毕设材料，准备答辩 \\
        \bottomrule
    \end{tabular}

    \vspace{1.5ex}
    \textbf{任务要求：}学生需要具备扎实的编程能力，熟练掌握 Python 语言，了解科学计算库如 NumPy、SciPy、Pandas 等的使用，有调用外部 API 和工具集成的经验。需要了解大语言模型的基本原理和应用，熟悉 LangChain、AutoGPT 等智能体框架，能够设计和实现工具调用的 Prompt 工程。具备跨学科的知识背景，对至少一个科学领域（如生物、化学、物理等）有基础了解，能够理解科学问题的特点和需求。

    \vspace{1ex}
    \noindent 起讫日期 2025 年 6 月 1 日 至 2026 年 6 月 3 日
    \begin{flushright}
        \bfseries \zihao{-4}
            指导教师（签名） \underline{\multido{}{5}{\quad}} 职称 \underline{\multido{}{5}{\quad}}
    \end{flushright}

    \noindent 三、系或研究所审核意见：\\

    \mbox{} \vfill
    \signature{负责人（签名）}
    % comment the line above and uncomment the line below if you want to set a signature with a specific date.
    % \signaturewithdate{负责人（签名）}{1897}{5}{21}
}
